% ==========================================
% SECCIONES 6 Y 7: ALCANCE Y LIMITACIONES
% ==========================================

\section{ALCANCE DEL PROYECTO}

El alcance del presente proyecto integrador abarca el ciclo de vida completo de desarrollo de software para el Sistema de Gestión de Doctorados (SGD), culminando en la entrega de una demostración tecnológica interactiva (demo funcional) completamente operativa, respaldada por manuales técnicos exhaustivos para su posterior adopción por parte de la institución. Con la finalidad de desligar recursos operativos y evitar dependencias de hardware durante la etapa académica, el prototipo se despliega sobre plataformas de nube (Vercel y Render), operando a partir de un banco de pruebas de expedientes sintéticos con perfiles e identidades ficticias generados a partir de los patrones oficiales de la institución para resguardar la privacidad y evitar la presencia de datos de personas reales en la demo pública.

En la fase inicial de gestión de información, el alcance comprende la revisión analítica de los formatos de las hojas de cálculo de Microsoft Excel oficiales utilizadas por la Unidad. A partir de dicha estructura institucional, se formula un banco de datos sintético y normalizado compuesto por expedientes y perfiles ficticios de ejemplo que preservan la representatividad lógica de los trámites y reglas de devengamiento, garantizando un riesgo nulo de exposición de información personal al no involucrar docentes reales en la demo pública, en estricto apego a la Ley Orgánica de Protección de Datos Personales (LOPDP). En el ámbito de persistencia, se implementa un modelo relacional en \textbf{PostgreSQL} gestionado dinámicamente con \textbf{Prisma ORM}, asegurando integridad referencial y restricciones transaccionales ACID.

En el componente de lógica de negocio y arquitectura desacoplada, el proyecto contempla la construcción de una capa de servicios API REST Backend modular desarrollada en \textbf{Express.js}, provista de middlewares para la autenticación por tokens JWT, sanitización de entradas y control de acceso basado en roles (RBAC). A nivel de interfaz de usuario, se desarrollará un \textbf{Panel Administrativo} en \textbf{Next.js 16 con React 19}, repotenciado por \textbf{Bun}, \textbf{Zustand}, \textbf{TailwindCSS} y \textbf{Driver.js}, concebido para que el personal de la Unidad (Director, Analista de Procesos y Analista Técnico) gestione integralmente los programas de doctorado y los expedientes de los estudiantes sin depender de hojas de cálculo de Excel. Este panel centralizará el control de requisitos documentales, la administración de la agenda y calendario de citas reservadas por los estudiantes, la gestión de adendas y el cálculo de plazos de devengamiento bajo seguridad RBAC.

Paralelamente, el alcance integra el desarrollo de un \textbf{Portal Estudiantil} en \textbf{Astro y React} enfocado en una carga ultrarrápida, provisto de consulta autónoma y transparente del estado del expediente docente (permitiendo al estudiante conocer en todo momento qué documentos y requisitos tiene aprobados y cuáles tiene pendientes de entrega para evitar visitas presenciales innecesarias), un sistema de reservas y agendamiento de citas en línea en tiempo real para concertar reuniones con cualquier miembro de la Unidad cuando requiera atención personalizada, un módulo de gestión de solicitudes y tickets, emisión digital y descarga de certificados de estado de beca, y un motor de notificaciones preventivas y alertas automáticas. La transferencia de conocimiento se consolidará mediante dos centros de documentación interactiva generados con \textbf{Fumadocs} (Manual de Usuario y Manual Técnico). Asimismo, el alcance comprende la ejecución de pruebas unitarias en el Backend, pruebas de integración y validación de requerimientos funcionales con la Unidad, culminando con la entrega de manuales técnicos interactivos (Fumadocs) y lineamientos de despliegue que orienten a los equipos técnicos de la universidad para una eventual migración institucional.

\section{LIMITACIONES DEL PROYECTO}

Las limitaciones identificadas para el desarrollo, despliegue y evaluación de la demostración tecnológica del SGD corresponden a factores de infraestructura externa, alcance del dataset y marco normativo institucional:

En primer lugar, respecto al entorno de datos y conectividad institucional, el prototipo opera de manera controlada sobre un banco de datos de prueba sintéticos con registros y perfiles ficticios de ejemplo. Para proteger la privacidad de la comunidad docente conforme a la LOPDP y desvincular la fase de titulación de autorizaciones complejas de red interna o accesos de escritura a bases de producción de la DTICs, la interoperabilidad institucional se modela mediante contratos de interfaz desacoplados y servicios simulados (\textit{mock services}). Esta decisión permite validar el 100\% de los flujos de consumo de datos de filiación y autenticación sin exponer a personas reales en servidores externos, dejando los conectores listos para enlazarse a los endpoints productivos de la DTICs durante la fase de migración definitiva.

En segundo lugar, en el ámbito de la infraestructura tecnológica y hospedaje, la adopción de plataformas de nube gestionadas (Vercel para las interfaces web y Render para el Backend y base de datos) exime a la universidad de aprovisionar servidores físicos inmediatos o configurar redes privadas virtuales (VPN) durante la fase de prototipado. Como contrapartida técnica, el entorno de demostración queda sujeto a las políticas operativas de dichas plataformas proveedoras, tales como una latencia inicial por reactivación de instancias tras periodos de inactividad (\textit{cold starts}) en el servicio de aplicaciones de Render, así como cuotas de ancho de banda y concurrencia asignadas por los proveedores para entornos de demostración pública. Para mitigar esta condición operativa y evitar interrupciones en las tareas programadas, el motor de alertas de vencimiento se diseña con un esquema desacoplado mediante disparadores externos programados y verificación reactiva bajo demanda, garantizando la continuidad de las notificaciones sin depender de procesos demonio residentes permanentes en el servidor.

Por último, en la dimensión reglamentaria, las fórmulas de cómputo de adendas y periodos de devengamiento se parametrizarán estrictamente bajo el marco legal de la Ley Orgánica de Educación Superior (LOES) y el Reglamento de Régimen Académico de la UCE vigentes en 2026 \cite{loes}, \cite{reglamento_uce}. Eventuales modificaciones o reformas normativas futuras en el Estatuto Universitario requerirán actualizaciones o adaptaciones posteriores en los módulos de lógica de negocio del sistema.

Asimismo, en lo concerniente a la selección de herramientas, lenguajes y frameworks tecnológicos detallados en este plan, se aclara que estos constituyen una propuesta técnica referencial sujeta a flexibilidad operativa. Su empleo efectivo o eventual sustitución a futuro dependerá directamente de los requerimientos, directrices o solicitudes específicas que formule la Unidad de Gestión de Doctorados a medida que avance el proyecto.

\subsection{Justificación Técnica del Banco de Datos Sintéticos Ficticios Frente a la Información Oficial}

Con la finalidad de precisar el alcance de los datos empleados en la plataforma y evitar cualquier interpretación errónea sobre el entorno de prueba, se detallan los fundamentos técnicos, éticos y operativos acordados formalmente con la Unidad de Gestión de Doctorados:

En primer lugar, el modelado conceptual de la base de datos relacional (PostgreSQL), la definición de esquemas en Prisma ORM y la lógica de controladores en Express.js fueron diseñados tomando como fuente directa las hojas de cálculo de Microsoft Excel oficiales utilizadas en la gestión cotidiana de la Unidad. Este análisis exhaustivo permitió estructurar con exactitud técnica las reglas de negocio reales (categorías de becas, causales de suspensión, cómputo de prórrogas y fórmulas de devengamiento). Sin embargo, el examen de estos archivos institucionales evidenció la presencia de inconsistencias históricas, campos vacíos e información personal de alta confidencialidad perteneciente a docentes becarios reales en funciones (nombres y apellidos, números de cédula de ciudadanía, remuneraciones económicas mensuales y antecedentes contractuales reservados).

En segundo lugar, al tratarse de un proyecto integrador cuya entrega académica involucra el despliegue de una demostración tecnológica interactiva (demo funcional) accesible vía web en plataformas de nube para su evaluación, la carga de datos de personas reales representaría un riesgo inaceptable de filtración o divulgación no autorizada de información confidencial, vulnerando las garantías de privacidad consagradas en la Ley Orgánica de Protección de Datos Personales (LOPDP) \cite{lopdp_ecuador}. Por este motivo, se acordó con la Unidad que para fines de la demo y las sesiones de prueba, la base de datos no contendría información de ningún docente real existente, sino que se poblaría mediante un \textbf{banco de datos sintéticos con perfiles e identidades enteramente ficticias}. Estos registros ficticios de ejemplo simulan el comportamiento de datos reales, reproduciendo fielmente la complejidad y variedad de los casos institucionales sobre la base de los formatos oficiales, pero asegurando un riesgo cero de exposición para los miembros de la comunidad universitaria.

En tercer lugar, esta estrategia garantiza que la solución de software quede plenamente concebida para una transición inmediata al entorno de producción institucional. Dado que los esquemas de tablas, tipos de datos y validaciones de negocio fueron programados calcando la arquitectura de trabajo real de la Unidad, la universidad no requerirá realizar modificaciones estructurales ni rediseñar el código fuente para su puesta en marcha oficial: bastará con aprovisionar la persistencia institucional y ejecutar la carga de los registros oficiales de los docentes becarios reales en la base de datos, siguiendo las directrices técnicas descritas en la \textbf{Guía Técnica para Desarrolladores} compilada en Fumadocs.

\subsection{Orientación Técnica y Transferibilidad para la Migración Institucional}

Para facilitar la adopción futura de la plataforma y evitar que la Unidad dependa de hojas de cálculo a largo plazo, el proyecto contempla la transferencia de conocimiento técnico mediante dos centros de documentación interactivos desarrollados en Fumadocs: el \textbf{Manual de Usuario} (dirigido al personal de la Unidad para la operación diaria) y la \textbf{Guía Técnica para Desarrolladores} (que documenta la arquitectura de software, los esquemas de bases de datos relacionales, las variables de entorno y los endpoints de la API REST). 

Es importante precisar que la ejecución técnica de una eventual migración de la base de datos a los servidores de producción de la Dirección de Tecnologías de la Información y Comunicación (DTICs), así como las gestiones de infraestructura de red interna y aprovisionamiento de servidores institucionales, corresponden a decisiones administrativas y contractuales propias de la universidad. No obstante, en concordancia con la cláusula de autoría intelectual y licenciamiento, los desarrolladores manifiestan su total predisposición para brindar acompañamiento técnico, asesoría y orientación a los equipos universitarios encargados de dicho proceso cuando sea formalmente requerido, aprovechando su conocimiento integral de la arquitectura desarrollada.